\section{構造方程式モデリング}
\subsection{授業内容}
\subsubsection{科目の中でのこのコマの位置づけ}
構造方程式モデリングは，因子分析と回帰分析を統合して扱う，総合的分析モデルである。言い換えれば，これまでの多くの多変量解析モデルのほとんどは，構造方程式モデルの下位モデルとして表現することができる。ここではこれまでのモデルを統合した，より現代的でより上位のモデルである構造方程式モデリングを理解することで，全ての多変量解析を網羅的かつ俯瞰的に捉えることが狙いである。

構造方程式モデリングを理解するには，変数の種類と関係性の区分に注意したパスダイアグラムの描き方を知ることが早い。パスダイアグラムを用いると，回帰分析と因子分析は説明変数が観測変数なのか潜在変数なのかといった違いであることが明らかである。また因子分析と似た主成分分析がどのように表されるかも，パスダイアグラムを見れば一目瞭然である。

パスダイアグラムには変数の尺度水準までかきこまれることはないが，ここに注意していろいろなモデルを描画すると，それがかつて多変量解析において様々な名称で呼ばれた分析方法であったことがわかる。あるいは，今後どのようなモデルが開発される可能性があるか，どのようなモデルをどのように希釈すれば良いかもイメージできる。

加えてこの統合的なモデルがなぜそうした複雑なモデルを表現できるのかについても，モデルを方程式で描画し，行列で考えることで，モデル行列とデータ行列を近づけることと理解することができる。この観点から，データにモデルを当てはめる適合度の考え方が改めて理解されるだろう。
\subsubsection{コマ主題細目}
\begin{description}
    \item[パスダイアグラムの書き方] これまで学んできたモデルを図で表現することを学ぶ。そのためには，変数を観測変数と潜在変数に区別することと，変数間関係を因果関係と相関関係に区別する必要がある。観測変数を矩形，潜在変数を楕円形，因果関係を一方向矢印，相関関係を双方向矢印で表現することで，因子分析や回帰分析が図で表現できることを学ぶ。

    \begin{flushright}
        $\to$\citet{小杉考司2014-01-10}のXXXXX.
    \end{flushright}


    \item[パスダイアグラムによる様々なモデル] 因子分析と回帰分析をパスダイアグラムで表現したことで，この両者を統合するような表現ができること，また潜在変数同士の関係を記述する，構造方程式を描画できるようになったと言える。因子分析と似た手法とされる主成分分析や，尺度水準の違いによる様々な統計モデルを表現する方法を手に入れたことになる。この手法を総称しt，構造方程式モデリングと呼ぶ。

    \begin{flushright}
        $\to$\citet{小杉考司2014-01-10}のXXXXX.
    \end{flushright}

    \item[構造方程式モデルによる未知数の推定] 構造方程式モデリングでは，パスダイアグラムでも表現されるが，変数間関係を方程式で書くこともできる。方程式で描画することで，構造方程式も潜在変数の方程式と観測変数の方程式，それらが入れ子になった方程式で描画できることがわかる。またこれらのモデルを行列のイメージで捉えると，最終的には分散共分散行列という実態を持った数字に対して，未知数で描画された方程式を接続したことが直感的にわかるだろう。未知数の増え方と分散共分散行列の要素の増え方を比較すると後者が圧倒的に早く，未知数よりも既知数が多い方程式は解くことができるという原理から，未知数が推定しうることを理解する。

    \begin{flushright}
        $\to$\citet{kosugi2018}のPp.191--193.
    \end{flushright}

    \item[適合度によるモデルの評価]  データ行列とモデル行列をイコールで結んだ方程式を解くことが，未知数を求める根本的な原理であるが，このことからモデルがデータとどの程度合致しているかという適合度が，モデル評価の統合的観点として浮かび上がってくる。回帰分析では$R^2$であったが，因子分析をはじめとした様々な多変量解析モデルも，この評価次元で考えることができる。ただしその指標にはいくつかの特徴があり，これらを総合的にみて評価するという実践的ノウハウも確認する。
\end{description}
\subsubsection{キーワード}
\begin{itemize}
    \item パスダイアグラム
    \item 観測変数と潜在変数
    \item 因果関係と相関関係
    \item 構造方程式モデリング
    \item 適合度
    \end{itemize}
\subsection{授業情報}
\paragraph{コマの展開方法}
講義
\subsubsection{予習・復習課題}
\paragraph{予習} 回帰分析と因子分析という二つの分析方法についてはすでに学んでいるが，この両者の共通点と相違点がどこにあるかを事前に考えてみよう。外的な基準の有無，説明変数の種類の観点から，自分の言葉で表現できるようになっていると良い。
\paragraph{復習} これまで学んだ様々な統計モデルを，構造方程式モデリングの表記法に則ってパス図を書いてみよう。また様々な尺度水準の組み合わせからなるモデルを考え，それらがどのような意味を持つのかと推論するのも理解の助けになる。